\Needspace{6\baselineskip}\SourceLeaf{86}

\[
= X^3\frac{d^2\varphi x}{dt\,du} + 3X^2\frac{dX}{dt}\frac{d\varphi x}{dt}
= \frac{dX^3 \frac{d\varphi x}{dt}}{dt},
\]

\noindent
so folgt

\begin{equation}\tag{5}
\frac{d^3\varphi x}{du^3} = \frac{d^2X^3 \frac{d\varphi x}{dt}}{dt^2}.
\end{equation}

\noindent\hspace*{\parindent}Ebenso findet man

\[
\frac{d^4\varphi x}{du^4} = \frac{d^3X^4 \frac{d\varphi x}{dt}}{dt^3} \quad \text{u.\,s.\,f.}
\]

\Needspace{9\baselineskip}\noindent\hspace*{\parindent}\textls{Allgemein aber} ergibt die Ableitung von $X^n\dfrac{d\varphi x}{dt}$ nach $u$ durch die
Substitutionen (3) und (4)

\begin{equation}\tag{6}
\frac{dX^n \frac{d\varphi x}{dt}}{du} = \frac{dX^{n+1} \frac{d\varphi x}{dt}}{dt}.
\end{equation}

\noindent
Gilt das Gesetz also bis zum Glied $n^{\text{te}}$, das heißt,

\[
\frac{d^n\varphi x}{du^n} = \frac{d^{n-1}X^n \frac{d\varphi x}{dt}}{dt^{n-1}},
\]

\noindent
so ergibt sich durch Differenzieren nach $u$

\[
\frac{d^{n+1}\varphi x}{du^{n+1}} = \frac{d\,d^{n-1}X^n\frac{d\varphi x}{dt}}{du . dt^{n-1}}
= \frac{d^{n-1}dX^n\frac{d\varphi x}{dt}}{dt^{n-1}du},
\]

\noindent
[also] nach (6)

\begin{equation}\tag{7}
\frac{d^{n+1}\varphi x}{du^{n+1}} = \frac{d^nX^{n+1} \frac{d\varphi x}{dt}}{dt^n},
\end{equation}

\noindent
d.\,h.\ [die Regel gilt] auch für das nächste Glied.

\noindent\hspace*{\parindent}Aus den Gleichungen (1), (4), (5) folgen also die obigen Umformungen für die ersten drei
Glieder, und aus (6) folgt ihre Gültigkeit für alle weiteren.

\medskip

\Needspace{10\baselineskip}\begin{center}
\textbf{BEMERKUNG ZUM AUFSATZ »NEUER BEWEIS DES LAGRANGESCHEN LEHRSATZES«.}
\end{center}

\noindent\hspace*{\parindent}Laut einer Notiz in dem auf S. 20 erwähnten Tagebuch hat \textsc{Gauss} am 27. Dezember 1796 einen
Beweis des Lehrsatzes von \textsc{Lagrange} gefunden. Der oben abgedruckte Aufsatz ist
zwar ohne Datumsangabe auf einem einzelnen Zettel aufgezeichnet. Es besteht
jedoch kein Zweifel, dass es sich hier um die Herleitung handelt, die \textsc{Gauss} Ende 1796 fand.
Die Veröffentlichung wurde damals nur durch verschiedene Zufälle verhindert
(vgl.\ \textsc{Sartorius von Waltershausen}, Gauss zum Gedächtniss, S. 22).

\begin{flushright}
\textsc{Fricke}.
\end{flushright}
